%------------------------------------------------------------------------------%
\documentclass[fleqn,utf8,aspectratio=169,12pt]{beamer}

\usepackage{gaal}

\title{Representação Matricial}

%------------------------------------------------------------------------------%
\begin{document}

\addtitle

%------------------------------------------------------------------------------%
\section{Representação Matricial}

%------------------------------------------------------------------------------%
\begin{frame}{Vetores como Matrizes Coluna}
  Podemos escrever um vetor de $\R^n$ como uma matriz coluna
  \[
    v = \begin{pmatrix}[c]
      v_1    \\
      v_2    \\
      \vdots \\
      v_n
    \end{pmatrix}
  \]
\end{frame}

%------------------------------------------------------------------------------%
\begin{frame}{Operações}
  \[
    v + w = \begin{pmatrix}[c]
      v_1+w_1 \\
      v_2+w_2 \\
      \vdots  \\
      v_n+w_n
    \end{pmatrix}
  \qquad\qquad\qquad\pause
    \alpha v = \begin{pmatrix}[l]
      \alpha v_1 \\
      \alpha v_2 \\
      \vdots     \\
      \alpha v_n
    \end{pmatrix}
  \]
\end{frame}

%------------------------------------------------------------------------------%
\begin{frame}{Produto Escalar}
  O produto escalar entre $v$ e $w$ é
  \pause
  \[
    v^t w
    \pause
    \;=\;
    \begin{pmatrix}
      v_1 & v_2 & \cdots & v_n
    \end{pmatrix}
    \begin{pmatrix}[c]
      w_1    \\
      w_2    \\
      \vdots \\
      w_n
    \end{pmatrix}
    \pause
    \;=\;
    \sum_{i=1}^{n} v_i w_i
  \]
\end{frame}

%------------------------------------------------------------------------------%
\section{Transformações Lineares}

%------------------------------------------------------------------------------%
\begin{frame}{Transformação linear}
  Uma aplicação (função)
  \[
    T\colon \R^n \to \R^m
  \]
  é linear quando
  \[
    T(v+w) = T(v) + T(w)
  \]
  \[
    T(cv) = cT(v)
  \]
\end{frame}

%------------------------------------------------------------------------------%
\begin{frame}{Combinações Lineares}
  As duas propriedades equivalem a
  \[
    T\left( c_1\vec{v}_1 + \cdots + c_k\vec{v}_k \right)
    = c_1 T(\vec{v}_1) + \cdots + c_kT(\vec{v}_k)
  \]
\end{frame}

%------------------------------------------------------------------------------%
%------------------------------------------------------------------------------%
\begin{question}
  Avalie a transformação linear
  \[
    T(x, y) = ( 2x-y, x+3y )
  \]
  em
  \[
    \vetor{x}{y} = \vetor{2}{3}
  \]
\end{question}

%------------------------------------------------------------------------------%
\begin{resolution}{Solução}
  \begin{align*}
    \onslide<+->{T(2, 3)}
    \onslide<+->{& = \vetor{2x-y}{x+3y}\evalat{x=2, \; y=3}{} \\}
    \onslide<+->{& = \vetor{2\times 2 - 3}{2 + 3 \times 3} \\}
    \onslide<+->{& = \vetor{1}{11}}
  \end{align*}
\end{resolution}

%------------------------------------------------------------------------------%
\endinput

%------------------------------------------------------------------------------%
\begin{question}
  Sabendo que
  \[
    T\left(\vec{i}\right) = \vetor{3}{5}
    \qquad\qquad
    T\left(\vec{j}\right) = \vetor{2}{1}
  \]
  calcule a transformação linear $T(\vec{v})$ para
  \[
    \vec{v} = \vetor{3}{4}
  \]
\end{question}

%------------------------------------------------------------------------------%
\begin{resolution}{Solução}
\begin{columns}[t]
\column{0.3\textwidth}
  \begin{align*}
    \onslide<+->{\vec{v}
    & = \vetor{3}{4} \\}
    \onslide<+->{& = \vetor{3}{0} + \vetor{0}{4} \\}
    \onslide<+->{& = 3 \vec{i} + 4 \vec{j}}
  \end{align*}
\column{0.6\textwidth}
  \begin{align*}
    \onslide<+->{T(\vec{v})}
    \onslide<+->{& = T\left(3 \vec{i} + 4 \vec{j}\right) \\}
    \onslide<+->{& = 3T\left(\vec{i}\right) + 4T\left(\vec{j}\right) \\}
    \onslide<+->{& = 3 \vetor{3}{5} + 4 \vetor{2}{1} \\}
    \onslide<+->{& = \vetor{9}{15} + \vetor{8}{4}}
    \onslide<+->{\;=\; \vetor{17}{19}}
  \end{align*}
\end{columns}
\end{resolution}

%------------------------------------------------------------------------------%
\endinput

%------------------------------------------------------------------------------%
\begin{question}
  Verifique que a transformação
  \(
    T\left(\vetor{x}{y}\right) = \vetor{2x-y}{x+3y}
  \)
  é linear
\end{question}

%------------------------------------------------------------------------------%
\begin{resolution}{Solução}
  Temos que mostrar que para quais quer vetores
  \[
    \vec{v} = \vetor{r}{s}
    \qquad\qquad
    \vec{w} = \vetor{p}{q}
  \]
  e qualquer escalar $c$

  \[
    T(\vec{v}+\vec{w}) = T(\vec{v}) + T(\vec{w})
  \]
  \[
    T(c\vec{v}) = c T(\vec{v})
  \]
\end{resolution}

%------------------------------------------------------------------------------%
\begin{resolution}{Soma}
  \begin{align*}
    \onslide<+->{T(\vec{v}+\vec{w})}
    \onslide<+->{  = T\left(\vetor{r}{s}+\vetor{p}{q}\right)}
    \onslide<+->{& = T\left(\vetor{r+p}{s+q}\right) \\}
    \onslide<+->{& = \vetor{2x-y}{x+3y} \evalat{x=r+p, \; y=s+q}{} \\}
    \onslide<+->{& = \vetor{2(r+p)-(s+q)}{(r+p)+3(s+q)} \\}
    \onslide<+->{& = \vetor{2p - q + 2r - s}{p + 3q + r + 3s}}
  \end{align*}
\end{resolution}

%------------------------------------------------------------------------------%
\begin{resolution}{Soma}
  \begin{align*}
    \onslide<+->{T(\vec{v}) + T(\vec{w})}
    \onslide<+->{& = T\left(\vetor{r}{s}\right) + T\left(\vetor{p}{q}\right) \\}
    \onslide<+->{& = \vetor{2x-y}{x+3y} \evalat{x=r, \; y=s}{}}
    \onslide<+->{  + \vetor{2x-y}{x+3y} \evalat{x=p, \; y=q}{} \\}
    \onslide<+->{& = \vetor{2r-s}{r+3s} }
    \onslide<+->{  + \vetor{2p-q}{p+3q} \\}
    \onslide<+->{& = \vetor{2p - q + 2r - s}{p + 3q + r + 3s} \\}
    \onslide<+->{& = T(\vec{v}+\vec{w})}
  \end{align*}
\end{resolution}

%------------------------------------------------------------------------------%
\begin{resolution}{Produto por Escalar}
    \begin{align*}
      \onslide<+->{T(c\vec{v})}
      \onslide<+->{  = T\left(\vetor{cr}{cs}\right)}
      \onslide<+->{& = \vetor{2x-y}{x+3y} \evalat{x=cr, \; y=cs}{} \\}
      \onslide<+->{& = \vetor{2cr-cs}{cr+3cs} \\}
      \onslide<+->{& = \vetor{c(2r-s)}{c(r+3s)} \\}
      \onslide<+->{& = c\vetor{2r-s}{r+3s}}
      \onslide<+->{\;=\; c T(\vec{v})}
    \end{align*}
\end{resolution}

%------------------------------------------------------------------------------%
\endinput
\end{frame}


%------------------------------------------------------------------------------%
\begin{frame}{Transformações Lineares e Matrizes}
  Uma transformação linear
  \[
    T\colon \R^n \to \R^m
  \]
  pode ser representada por uma matriz
  \[
    A \in \R^{m\times n}
  \]
  tal que
  \[
    T(v) = A v
  \]
\end{frame}

%------------------------------------------------------------------------------%
%------------------------------------------------------------------------------%
\begin{question}
  Determine a forma matricial \(T(v)=Av\) da transformação linear
  \[
    T(x, y) = \vetor{2x-y}{x+3y}
  \]
\end{question}

%------------------------------------------------------------------------------%
\begin{resolution}{Solução}
  Precisamos determinar uma matriz $A$
  \pause
  tal que
  \[
    A\vetor{x}{y} = \vetor{2x-y}{x+3y}
  \]

  \pause
  Como
  \[
    \begin{pmatrix}
      2 & -1 \\
      1 & 3
    \end{pmatrix}
    \begin{pmatrix}
      x \\
      y
    \end{pmatrix}
    \;=\;
    \begin{pmatrix}
      2x-y \\
      x+3y
    \end{pmatrix}
  \]

  \pause
  \[
    A = \begin{pmatrix}
          2 & -1 \\
          1 & 3
        \end{pmatrix}
  \]
\end{resolution}

%------------------------------------------------------------------------------%
\endinput


%%------------------------------------------------------------------------------%
%\section{Lista Mínima}
%
%%------------------------------------------------------------------------------%
%\begin{frame}{Lista Mínima}
%  \begin{enumerate}
%    \item
%  \end{enumerate}
%
%  \vfill
%  Atenção: A prova é baseada no livro, não nas apresentações
%\end{frame}

%------------------------------------------------------------------------------%
\end{document}
%------------------------------------------------------------------------------%
